% !PRINTED 91
\weberpage{91}

%<<CONT
auch in diesen F\"allen $e_{1,h}$ so annehmen, dass $c_{1,h}\,e_{1,h}$
durch $q_h$ nicht theilbar wird.

Die Forderungen, die hiermit f\"ur die $e_{k,h}$ gestellt sind, bestehen
also nur darin, dass sie durch gewisse Primzahlen theilbar, durch andere
nicht theilbar sein sollen, und sind also noch auf unendlich viele Arten
mit einander vertr\"aglich. Aber diese letzte Annahme \"uber die $e_{k,h}$
hatte nur den Zweck, zu zeigen, dass die fr\"uheren allgemeineren
Forderungen immer befriedigt werden k\"onnen. Es kann noch viele andere
Arten geben, ihnen zu gen\"ugen.

\noindent\rule{4cm}{0.4pt}

\begin{center}
\gesperrt{Uebersicht der Resultate.}
\end{center}

Wenn es sich darum handelt, alle Kreistheilungsk\"orper, und jeden nur
einmal, und zwar auf m\"oglichst einfache Art, durch
Kreistheilungsperioden darzustellen, so verfahre man so:

Man nehme ein beliebiges System von Primzahlpotenzen als Invarianten

\[
i_1, i_2, \dots i_v
\]

an, und w\"ahle dann eine Zahl $m$ nach folgenden Bedingungen.

Wenn $p^\pi$ die h\"ochste Potenz einer Primzahl $p$ ist, die unter den
Invarianten vorkommt, so nehme man $p$ h\"ochstens $(\pi+1)$-, oder,
wenn $p = 2$ ist, $(\pi+2)$ mal in $m$ auf.

Eine Primzahl $q$, die gar nicht in den Invarianten aufgeht, darf nur
einfach in $m$ aufgenommen werden; aber nur solche einfache Primfactoren
$q$ darf $m$ haben, bei denen $q - 1$ durch einen der Primfactoren $p$
der Invarianten theilbar ist. Also kann auch $p$ selbst nur dann einfach
in $m$ vorkommen, wenn $p - 1$ durch eines der anderen $p$ theilbar ist.

Ferner muss $m$ noch der Bedingung gen\"ugen, dass unter den
Indexmoduln

\[
c_1, c_2, \dots c_\mu
\]

$\varrho$ verschiedene, wie $c_1, c_2, \dots c_\varrho$, durch $i_1, i_2,
\dots i_\varrho$ theilbar sind, wenn $i_1, i_2, \dots i_\varrho$ Potenzen
von einer und derselben Primzahl sind. Dann bestimme man die Gr\"ossen
$i_{k,h}$ und $c_{k,h}$ nach den Formeln (11) und theile die gegebenen
Invarianten in Systeme ein, so dass alle Potenzen einer und derselben
Primzahl in einem Systeme vereinigt sind.

F\"ur jedes dieser Systeme bilde man die Matrix (24) und w\"ahle die
Zahlen $c_{k,h}$ so, dass aus jeder solchen Matrix eine
% [sic: Weber prints c_{k,h}; these quantities were fixed by (11), whereas the matrix entries and the ensuing construction require choosing e_{k,h}.]
%CONT>>
